\section{Model index and tool migration}\label{sec:C17:migration}
Memory portability depends on the consumer. A concise skill written for one model can rely on its implicit knowledge or compensate for its characteristic mistakes. A new model may interpret the same text differently. A new embedding model can change the geometry of queries and records. A new tool schema can invalidate a previously successful procedure. Version identifiers make these changes visible but do not establish compatibility.

An embedding example makes the issue precise. Suppose document vectors $d$ and queries $q$ originally share a space with score $q^{\mathsf T}d$. A replacement encoder applies an orthogonal rotation $R$ to both, preserving within-space dot products. If only queries are updated, the score becomes $(Rq)^{\mathsf T}d$, which generally differs from $q^{\mathsf T}d$. Even an information-preserving coordinate change breaks mixed-version retrieval. Re-embedding both sides or using a validated alignment is required. This is a mathematical example, not a benchmark result about named embeddings.

For policy migration, preserve the old memory snapshot and compare the new model on matched tasks with no memory, original memory, and a revised memory representation. This separates backbone changes from memory compatibility. Include cases where the old memory should be ignored. A new model that applies every old instruction more aggressively can improve some recall tasks while violating current constraints. Positive transfer and overuse need separate outcomes.

Tool migration can be more discrete. Changing an argument's default, resource scope, or side-effect semantics can invalidate a skill even when its function name remains unchanged. Attach schema and behavior versions to the skill's compatibility record, and require fresh validation when relevant contracts change. A schema-compatible change can still alter semantics; syntactic validation is necessary but insufficient.

CLM's inspected environment shows a useful local boundary: \texttt{ContextEnv.step} writes the live mirror, executes a command, reads the mirror back, and applies a parsed edit through \path{_apply_edit}. The edit gate checks growth against a configured rule, and the parser retains a protected prefix~\cite{repo_clm2026}. This validates some structural and budget conditions, not the truth of every rewritten assertion. Migrating that environment to a new tokenizer or parser changes what the gate counts and which messages it reconstructs.

A successful edit is also not synonymous with fitting the token budget. In the inspected \path{_apply_edit} branch, growth is subject to the configured growth/fit rule, but a shrinking edit can be applied while the result remains over the limit; the readout then requests further compaction~\cite{repo_clm2026}. The synchronization counter likewise does not prove behavioral preservation. An integration should record applied versus rejected edits, final enforced token counts, and the actual next-decision input separately. Structural acceptance supplies no bound on $\mathsf P_\epsilon$.

A migration gate should therefore cover the entire path from stored record to dispatched effect. Record old and new versions, representative compatibility tests, unresolved failures, and rollback criteria. If a proprietary model version cannot be pinned exactly, state that reproducibility limit and use dated API identifiers and observed behavior. A rolling service name is not a reproducible policy identity.

The October 2026 skill--policy co-evolution study makes this dependence experimentally explicit through CODA. It perturbs the delivered skill bank and evaluates consistency under meaning-preserving rewrites, resistance to corrupted guidance, and worst-case success~\cite{feng2026coda}. Its theory concerns candidate coverage and a specified cooperative acceptance rule, not every conceivable self-improvement objective. Its stress-test losses also do not estimate how frequently equivalent drift occurs in deployment. The result is a concrete warning against evaluating only the policy-bank pair that was selected together.

A corresponding original diagnosis separates three perturbations. Rewrite surface form while preserving procedural meaning, alter valid preconditions while retaining familiar wording, and replace the consumer policy while freezing the bank. The first tests interpretive invariance, the second tests applicability sensitivity, and the third tests portability. Pooling them into one robustness number can reward the wrong behavior: a policy should ignore irrelevant wording changes but respond to a changed precondition. Consistency is therefore invariance to specified transformations and sensitivity to others, not indiscriminate resistance to change.

